\section{Combining the logical encodings: a complete 692-bit packet}
\label{sec:eq704}

There is a successful combination of two genuinely different ideas in the
inspected intake, even though it is not by itself a binary-state reduction.
The eight-equation packet $\mathcal E_{713}$ uses a complemented existential
binder and a factored closure equation.  The separate 937-bit packet uses
three structural constructors as closure witnesses.  Applying those three
identifications to $\mathcal E_{713}$ gives a complete \textbf{704-bit}
equational consistency problem: nine operation symbols, eight equations,
159 equation-term nodes, four local variables, and two ground endpoints.
That first packet is exactly 88 bytes.  The nine saved bits are all
signature bits; the term-node count is unchanged.  A further Boolean
identity reduces its closure equation from 82 to 80 term nodes, producing
the \textbf{696-bit checkpoint $\mathcal E_{696}$}, exactly 87 bytes.
A further structural identification uses the already unary $Q$ as
selector successor, removing one more term node and giving the
\textbf{692-bit packet $\mathcal E_{692}$}.  Its packed representation
uses 87 bytes with four trailing padding bits.  We retain both checkpoints
so that the model combination, Boolean identity, and change of selector
representation can be assessed separately.

The substantive issue is compatibility of the model interpretations.
In the new packet $Q$ must mean existential negation as a formula
constructor and must also supply a closure hull or a bad input as a raw
object.  We prove that these roles coexist by a well-founded assignment
on a free tagged carrier.  In particular, reusing the 937-bit packet's
ordinary-existential interpretation of $Q$ without changing it would be
incorrect.  The first new relative theorem is
\begin{equation}
 \mathrm{ZFC}\vdash
 \bigl(\mathcal E_{704}\nvdash N(d,d)=d
       \ \longleftrightarrow\
       \mathcal E_{713}\nvdash N(d,d)=d\bigr).
 \label{eq:704-relative}
\end{equation}
Its proof below needs no new assumption about the consistency strength
of raw EPSC.  Composing it with the inherited ordinary calibration of
$\mathcal E_{713}$ gives the same equivalence with
$\operatorname{Con}(\mathrm{ZFC})$.  This is an ordinary mathematical
proof formalizable in ZFC.  A serialized Lean/Rocq endpoint and an
accounted binary interpreter for the packet are not supplied here.

\subsection{Exact equations and encoding}

The ordered signature is
\[
 d^0,N^2,V^2,P^2,A^2,Q^1,W^2,p_E^2,p_\varepsilon^2.
\]
The first seven equations are
\begin{align}
 N(N(x,N(N(y,x),x)),N(y,N(z,x)))&=y,\label{eq:704-sh1}\\
 d&=N(d,N(d,d)),\label{eq:704-top}\\
 V(P(x,s),d)&=x,\label{eq:704-head}\\
 V(P(x,s),A(j,d))&=V(s,j),\label{eq:704-tail}\\
 p_\varepsilon(s,P(a,b))
   &=N(p_\varepsilon(s,a),p_\varepsilon(s,b)),\label{eq:704-nand}\\
 p_\varepsilon(s,A(i,j))
   &=p_E(V(s,i),V(s,j)),\label{eq:704-atom}\\
 p_\varepsilon(s,Q(a))
   &=N\bigl(p_\varepsilon(P(W(s,a),s),a),
            p_\varepsilon(P(z,s),a)\bigr).
       \label{eq:704-binder}
\end{align}
The first equation is the published Sh1 single axiom for Boolean algebra,
not a newly discovered axiom.  The second makes $d$ the Boolean top.
Equation~\eqref{eq:704-binder} is universal in $z$ as well as its other
variables.  As explained below, its intended binder is $\exists\neg$.

For the eighth equation the following are only abbreviations for terms:
\begin{align*}
 h&=Q(b),&n_z&=P(b,z),&m_z&=A(b,z),\\
 R_b(x,y)&=p_\varepsilon(P(y,P(x,b)),V(b,d)),&&&
\end{align*}
and
\begin{align*}
 r&=R_b(h,n_z),&e&=p_E(m_z,n_z),&f&=p_E(m_z,z),\\
 a&=p_E(W(x,d),x),&b_0&=p_E(x,h),&c&=R_b(x,y),\\
 u&=p_E(y,h),&v&=R_b(P(b,y),y).&&
\end{align*}
Put
\[
 \ell=N(d,N(c,N(a,N(d,b_0)))),\qquad
 B^-=N(r,N(d,N(e,N(d,f)))),
\]
\[
 H^-=N\bigl(N(\ell,N(d,u)),N(u,N(d,v))\bigr).
\]
The final equation is
\begin{equation}
                         d=N(B^-,H^-).
 \label{eq:704-closure}
\end{equation}
The files \texttt{E704.equations}, \texttt{E704.json}, and
\texttt{E704.bits} contain the fully expanded terms.  No abbreviation
or signature information is supplied outside the bit packet.

This is the literal substitution
\begin{equation}
 \mathsf{Target}(b)=Q(b),\qquad
 \mathsf{Next}(b,z)=P(b,z),\qquad
 \mathsf{Mem}(b,z)=A(b,z)
 \label{eq:704-aliases}
\end{equation}
in $\mathcal E_{713}$.  In particular, every separating
$\mathcal E_{704}$-model expands to a separating
$\mathcal E_{713}$-model by these definitions.  That direction is
immediate; the converse requires the following construction.

\subsection{What the inherited packet implies about raw membership}

We give enough detail to make the new relative theorem independent of
a black-box assertion about raw EPSC.  Work with a nonempty domain and
a binary relation $E$, with no equality in the formulas of the schema.
Write
\begin{align*}
 \mathsf{Empty}(x)&\iff\forall t\,\neg E(t,x),\\
 y\subseteq_E z&\iff\forall t\,(E(t,y)\to E(t,z)),\\
 \mathsf{LB}(R)&\iff
       \forall x\,\exists z\,\forall y\,(R(x,y)\to y\subseteq_E z).
\end{align*}
For every finite pure-membership formula $R(x,y;\bar p)$ and every
finite parameter tuple, raw EPSC asserts
\begin{equation}
 \mathsf{LB}(R)\ \longrightarrow\
 \exists h\left[
   \begin{aligned}
   &\forall x,y\,
     (((\mathsf{Empty}(x)\lor E(x,h))\land R(x,y))\to E(y,h))\\[-1mm]
   &{}\land\forall y\,(E(y,h)\to\exists x\,R(x,y))
   \end{aligned}\right].
 \label{eq:704-epsc}
\end{equation}
This is a syntactic first-order schema.  It is not collection for all
semantic binary relations.

\paragraph{Complemented-witness fact.}
For a Boolean-valued family $(u_z)$ and a selected value $w=u_{z_0}$,
\[
 \forall z\,[q=\neg(w\wedge u_z)]
 \quad\Longleftrightarrow\quad
 q=\neg w\ \land\ \forall z\,[w\leq u_z].
\]
Taking $z=z_0$ proves the first conjunct; complementing the remaining
equations proves the second.  After a Boolean homomorphism to two
values this says $q\iff\exists z\,\neg u_z$, with the selected
instance false whenever such an instance exists.  It does not assert
the existence of an attained minimum in an arbitrary Boolean-valued
family.

\paragraph{Extraction lemma.}
Every countable separating $\mathcal E_{713}$-model gives a countable
raw-EPSC structure on the same nonempty carrier.

\begin{proof}
Sh1 makes the carrier a nontrivial Boolean algebra.  Choose an
ultrafilter and its homomorphism to two truth values.  Apply it to
the values of $p_E$ and $p_\varepsilon$; retain all native functions
and their literal equality on the original carrier.  The complemented
witness fact proves the $\exists\neg$ semantics of $Q$ for the resulting
class relation $\varepsilon$.

Let $\nu_0=d$ and $\nu_{i+1}=A(\nu_i,d)$.  The two native stack
equations recover the $i$th coordinate of a finite stack by repeated
tail equations followed by one head equation.  Atomic membership
is represented by $A(\nu_i,\nu_j)$, NAND by $P$, and ordinary
existential quantification of a class $a$ by $Q(P(a,a))$.
Structural induction, shifting indices under binders, represents every
finite equality-free pure-membership formula.  For a matrix in $x,y$
and parameters, use output coordinate 0, input coordinate 1, skip
coordinate 2, and put parameters at coordinates 3 onward.  Build its
class code $c$ first and then the descriptor
$b=P(c,P(p_1,\ldots P(p_k,d)\ldots))$.  Its represented relation
is exactly $R_b$.  This is a finite term construction without circular
definition.  Logical truth and falsity, if needed, are represented by
NAND tautologies and contradictions built from an atomic class; the
lower argument does not assume that the arbitrary class $d$ is false.

In two truth values, the final equation of $\mathcal E_{713}$ is
\[
 \forall b,z,x,y\quad
 \mathsf{Bad}(b,z)\lor\mathsf{Hull}(b,x,y),
\]
where
\begin{align*}
 \mathsf{Bad}&=r\land e\land\neg f,\\
 \mathsf{Hull}&=(a\lor\neg c\lor u)
            \land(\neg b_0\lor\neg c\lor u)
            \land(\neg u\lor v).
\end{align*}
For a represented $R_b$ satisfying $\mathsf{LB}$, apply that hypothesis
at $h=\mathsf{Target}(b)$.  It supplies one $z$ for which the three
bad literals cannot all hold.  The hull clauses therefore hold for
every $x,y$.  If $x$ is empty, $E(W(x,d),x)$ is false, so the first
hull clause gives closure at empty inputs.  The second gives closure
at members of $h$.  The third gives unrestricted support, with witness
$\mathsf{Next}(b,y)$.  These are precisely~\eqref{eq:704-epsc}.
\end{proof}

\subsection{A model of the combined packet from any countable EPSC structure}

Let $M=(M,E_M)$ be a countable nonempty raw-EPSC structure, and fix
$m_0\in M$.  Take the carrier $D$ to consist of all finite free trees
whose leaves are a distinguished $d$ and separate elements $e_m$ for
$m\in M$, with constructors $P,A,Q$.  The leaves $e_m$ belong to the
chosen model carrier; they are not extra constant symbols in the
finite packet.  Interpret the three constructors as their literal
disjoint injective tree constructors.

Give leaves weight 1 and put
\begin{equation}
 w(Q(b))=w(b)+1,\quad
 w(P(b,z))=w(b)+w(z)+2,\quad
 w(A(b,z))=w(b)+w(z)+3.
 \label{eq:704-weight}
\end{equation}
Every proper subtree has smaller weight, and
\[
               w(Q(b))<w(P(b,z))<w(A(b,z)).
\]
These inequalities order the hull choice, output witness, and separating
member witness.  They are the reason the structural identifications
are compatible.

For any selector $i$, define $\operatorname{idx}(i)$ to be the number
of consecutive outer layers $A(j,d)$; a selector of any other shape
has index zero.  Read $P(x,s)$ as a cons cell and every non-$P$ tail
as the infinite constant-$d$ stream.  Define $V(s,i)$ to return
coordinate $\operatorname{idx}(i)$.  The two native stack equations
hold for all carrier arguments, including noncanonical selectors.

Decode a finite class tree structurally: each leaf is falsity, $P$
is NAND, $A(i,j)$ is membership between coordinates
$\operatorname{idx}(i),\operatorname{idx}(j)$, and $Q(a)$ prepends
one coordinate and existentially quantifies the negation of its body.
Thus $Q$ is interpreted as $\exists\neg$, not as the ordinary
existential of the 937-bit input.  This decoding is pure syntax and
does not use a raw-value assignment.

\paragraph{Descriptor descent.}
For every $b$, the class expression
$\varepsilon(P(y,P(x,b)),V(b,d))$ describes a finite pure-membership
formula $\widehat R_b(x,y;\bar u)$.  Every fixed parameter in
$\bar u$ is a proper subtree of $b$ or the default leaf $d$.
Indeed, the class $V(b,d)$ is the head of $b$ when $b$ is a $P$-tree
and is $d$ otherwise.  Its finitely many referenced coordinates are
the two new inputs, bound coordinates, or coordinates of the tail
$b$.  Each of the latter is the head of a proper tail subtree or the
default leaf.  The argument also covers codes referring to their own
code slot and arbitrarily deep selector indices.  Changing an
existential binder to $\exists\neg$ does not change this descent.

Define a projection $\mu:D\to M$ by recursion on weight, starting
with $\mu(d)=m_0$ and $\mu(e_m)=m$.  For each descriptor $b$, after
the proper-subtree parameters of $\widehat R_b$ have received raw
values, let $R_b^M$ be the resulting pure-membership relation in $M$.
Use a fixed enumeration of $M$ to make the following choices consistently
whenever the same descriptor recurs.

If $\mathsf{LB}(R_b^M)$ holds, choose an EPSC hull $h_b$ and mark
this the hull case.  Otherwise choose $h_b$ satisfying
\begin{equation}
 \forall z\in M\,\exists n,m\in M\quad
 R_b^M(h_b,n)\land E_M(m,n)\land\neg E_M(m,z).
 \label{eq:704-bad-input}
\end{equation}
This is exactly classical negation of local boundedness.  Set
$\mu(Q(b))=h_b$.

In the bad-input case, choose $\mu(P(b,z))=n$ with
$R_b^M(h_b,n)$ and $n\not\subseteq_{E_M}\mu(z)$; then choose
$\mu(A(b,z))$ to be a member of $n$ outside $\mu(z)$.
In the hull case, if $E_M(\mu(z),h_b)$, choose
$\mu(P(b,z))=n$ with $R_b^M(n,\mu(z))$, using hull support;
otherwise use $m_0$.  In this case put $\mu(A(b,z))=m_0$.

This is a total recursion.  At $Q(b)$ all descriptor parameters have
smaller weight.  At $P(b,z)$, the value of $Q(b)$ and that of $z$
are available.  At $A(b,z)$, the value of $P(b,z)$ is available.
Every requested witness exists in its selected case.  No truth
predicate for the expanded language and no circular definition of
formula meaning is used.  The leaves $e_m$ make $\mu$ surjective.

Define $E_D(u,v)$ by $E_M(\mu(u),\mu(v))$.  For every
equality-free pure-membership formula $\phi$,
\begin{equation}
 (D,E_D)\models\phi(\bar u)
 \quad\Longleftrightarrow\quad
 (M,E_M)\models\phi(\mu(\bar u)).
 \label{eq:704-transfer}
\end{equation}
Induction on formulas proves this; surjectivity is used at quantifiers.
Literal equality is excluded from this transfer because distinct free
tags can have the same raw value.  Native equality in the equational
packet remains literal equality of trees throughout.

Let $\varepsilon(s,c)$ be the truth value of the finite class decoded
from $c$ on stream $s$ in $(D,E_D)$.  For $a\ne d$, choose $W(s,a)$
to falsify $\varepsilon(P(z,s),a)$ whenever such a $z$ exists,
and choose $d$ otherwise.  Countability permits a fixed least-index
choice.  At $a=d$, every class instance is false and the
$\exists\neg$ class $Q(d)$ is true.  Thus both sides of the binder
equation are true independently of $W(s,d)$, leaving that slice free.
Allocate it by
\[
 W(x,d)=
 \begin{cases}
 e_m\text{ for a fixed selected }m\text{ with }E_M(m,\mu(x)),
       &\text{if such }m\text{ exists},\\
 d,&\text{otherwise}.
 \end{cases}
\]
This makes $E_D(W(x,d),x)$ true whenever $x$ is nonempty.  There is
only one allocation to this free slice.

Formula transfer identifies the relation represented by $R_b$ with the pullback
\[
 (x,y)\longmapsto R_b^M(\mu(x),\mu(y)).
\]
In the bad-input case, the three raw values
\[
 \mu(Q(b)),\qquad\mu(P(b,z)),\qquad\mu(A(b,z))
\]
make every bad clause true.
In the hull case, the first hull clause follows from member choice at
nonempty $x$ and from EPSC closure at empty $x$; the second follows
from closure at members of the hull; the third follows from the
selected support witness $\mu(P(b,y))$.  Hence
$\mathsf{Bad}(b,z)\lor\mathsf{Hull}(b,x,y)$ holds for all tuples.
No Boolean branch flag is needed in the signature.

Finally transport a countably infinite Boolean algebra, for example
the finite/cofinite algebra, onto $D$ with top equal to its distinguished
element $d$.  Interpret $N$ as NAND, and $p_E,p_\varepsilon$ as that
algebra's top or bottom according to $E_D,\varepsilon$.  Keep every
native function fixed.  All eight equations now hold literally and
$N(d,d)\ne d$.  This completes the upper-model construction.

\subsection{Ground nonderivability, calibration, and the exact remaining cost}

For either finite equational packet, the presence of $d$ ensures a
nonempty countable algebra of ground terms.  Let $\sim$ be the least
congruence containing all ground instances of both orientations of all
equations.  If the endpoints are not equationally derivable, their
classes in the ground-term quotient are distinct.  That quotient is
a countable separating model: choose ground representatives for the
finitely many values assigned to the variables of an equation.
Conversely, equational soundness prevents a separating model from
equating the endpoints derivably.

Starting with a separating $\mathcal E_{704}$-model, expand by
the three literal definitions to obtain an $\mathcal E_{713}$-model.
Starting with nonderivability for $\mathcal E_{713}$, take its
countable ground-term model, apply the extraction lemma, and use
the weighted upper construction to obtain a separating
$\mathcal E_{704}$-model.  This proves~\eqref{eq:704-relative}.
All objects in this argument are sets; recursion is on finite syntax
or positive integer weights, and the uses of ultrafilters and ordinary
first-order soundness are legitimate in ZFC.

The inherited dependency for the full calibration is the ordinary
raw-EPSC/ZFC comparison in ProveIt's Sources 14 and 08, and its use in
the $\mathcal E_{713}$ argument.  With that dependency explicitly
retained, the conclusion is
\[
 \mathrm{ZFC}\vdash
 \bigl(\mathcal E_{704}\nvdash N(d,d)=d
             \ \longleftrightarrow\ \operatorname{Con}(\mathrm{ZFC})\bigr).
\]
Thus a fair enumerator of ground contextual rewrite certificates
from $N(d,d)$ to $d$ is a new precisely specified consistency search
target.  Every certificate must assign a ground value to the union
of variables on both sides of a chosen equation, including variables
absent from the matched side when the equation is reversed.

The unchanged codec uses an 8-bit symbol count, three arity bits per
symbol, an 8-bit variable capacity, one leading bit per equation,
one zero list terminator, and fixed-width term opcodes.  Nine symbols
plus four variables require four bits per term node.  The complete
cost is
\[
 \underbrace{8+3\cdot9+8}_{43}
 +\underbrace{8+1+4\cdot159}_{645}
 +\underbrace{4\cdot4}_{16}
 =\boxed{704\text{ bits}}.
\]
The equation-node counts are $14,6,6,10,12,12,17,82$.  The bit-character
SHA-256, excluding the final newline, is
\begin{center}\scriptsize\ttfamily
616d4eef0e70b67947b79ca5edc23dab718c06efd0c730a8e0929850b179629f.
\end{center}
The builder reconstructs the exact incoming 713-bit packet before
specializing it, roundtrips the result, rejects all 704 strict
prefixes and two malformed controls, and checks the closure Boolean
identity on all 256 assignments.  These checks validate syntax and
transcription; the unbounded model proof above supplies the logical
argument.  No finite validation is described as a model of EPSC.

\paragraph{Further research.}
The next comparison must charge a ground-rewrite verifier, substitution,
context navigation, complete enumeration, reset, and a blank-tape
initializer in the same binary TM convention as the existing candidate.
The 704 bits alone do not imply any state bound.  A particularly useful
next proof artifact would check the weighted descriptor recursion and
the two binder interpretations in one kernel development.  The
signature sharing may be reusable in other finite interpreted theories,
but each proposed sharing must exhibit its own well-founded dependency
order and preserve any unconstrained witness slice it uses.

\subsection{The 696-bit checkpoint: eliminating a double complement}
\label{sec:eq696}

The 704-bit checkpoint implements the hull requirement as the interval
condition $\ell\leq u\leq v$, where
$\ell=c\wedge(\neg a\lor b_0)$.  Its complement is used in the
outer closure equation.  The same condition admits the Boolean
identity
\begin{equation}
 (\neg\ell\lor u)\wedge(\neg u\lor v)
   =(\neg u\wedge\neg\ell)\lor(u\wedge v).
 \label{eq:696-mux}
\end{equation}
This identity holds in every Boolean algebra: distribute the left
side, then absorb its consensus term $\neg\ell\wedge v$ into the
other two terms, using $u\lor\neg u=1$.

The complement of $\ell$ has the cheaper direct expression
\[
                 K=N(c,N(a,N(d,b_0))).
\]
Define
\begin{equation}
 H=N\bigl(N(N(d,u),K),N(u,v)\bigr).
 \label{eq:696-hull}
\end{equation}
The NAND meaning of this term is the right side
of~\eqref{eq:696-mux}.  The 696-bit packet retains the first seven
equations of $\mathcal E_{704}$ and replaces the eighth equation by
\begin{equation}
                        d=N(B^-,N(d,H)),
 \label{eq:696-closure}
\end{equation}
where $B^-$ and all predicate-term abbreviations are unchanged.
Every abbreviation is expanded in \texttt{E696.equations} and the
encoded packet.

The previous expression first formed $\ell$ by complementing $K$;
the new expression uses $K$ directly.  After the complete NAND
expressions are expanded, this saves two term nodes in the closure
equation.  The remaining appearances of $u$, every relation witness,
the membership support clause, and both closure clauses are retained.
The change is a Boolean identity, not a weakening of EPSC.

\paragraph{Boolean equivalence theorem.}
$\mathcal E_{696}$ and $\mathcal E_{704}$ have exactly the same
models in their common signature, and hence exactly the same
equational consequences.  In particular, in the ordinary ZFC
metatheory,
\[
 \mathcal E_{696}\nvdash N(d,d)=d
 \quad\Longleftrightarrow\quad
 \mathcal E_{704}\nvdash N(d,d)=d
 \quad\Longleftrightarrow\quad
 \mathcal E_{713}\nvdash N(d,d)=d.
\]
To prove the first assertion, use Sh1 and the top equation to interpret
the common reduct as a Boolean algebra.  Identity~\eqref{eq:696-mux}
makes the two eighth equations equivalent for every assignment;
the other equations are identical.  Ground-term equational
completeness, or the earlier separating-model argument, gives the
assertion about the designated query.  Composing with the explicitly
retained E713 calibration yields
\[
 \mathrm{ZFC}\vdash
 \bigl(\mathcal E_{696}\nvdash N(d,d)=d
          \ \longleftrightarrow\ \operatorname{Con}(\mathrm{ZFC})\bigr).
\]

The 696-bit equation-node counts are
$14,6,6,10,12,12,17,80$, totaling 157.  Under the unchanged codec,
\[
 \underbrace{8+3\cdot9+8}_{43\text{ header bits}}
 +\underbrace{8+1+4\cdot157}_{637\text{ equation bits}}
 +\underbrace{4\cdot4}_{16\text{ endpoint bits}}
 =\boxed{696\text{ bits}}.
\]
Its bit-character SHA-256, excluding the final newline, is
\begin{center}\scriptsize\ttfamily
d0e84abc889413e5b12781b7aa3176901be6bd0a97cccbc40dc95b04ee35302a.
\end{center}
The builder checks the full closure identity and the new hull identity
on all 256 assignments to the eight independent Boolean predicates,
roundtrips both complete packets, rejects all 696 strict prefixes of
the 696-bit checkpoint, and rejects two additional malformed controls.  It
retains E704 as an independently encoded checkpoint and reconstructs
the exact incoming E713 baseline.  These are syntax and Boolean
checks; neither the 87-byte size nor the model-equivalence theorem
charges the missing binary rewrite interpreter or its initializer.

\subsection{The final 692-bit packet: sharing the unary binder with selectors}
\label{sec:eq692}

The remaining selector successor $A(j,d)$ costs three term nodes.
The signature already contains the unary constructor $Q$.  Replace
only the fourth equation of $\mathcal E_{696}$ by
\begin{equation}
                 V(P(x,s),Q(j))=V(s,j).
 \label{eq:692-tail}
\end{equation}
Keep all seven other equations, the ordered signature, and the
designated ground endpoints unchanged.  This defines
$\mathcal E_{692}$.  The change is a new representation of finite
indices, not a Boolean identity or a claim that the two packets
have identical models.

\paragraph{Selector-sharing theorem.}
In the ordinary ZFC metatheory,
\[
 \mathcal E_{692}\nvdash N(d,d)=d
 \quad\Longleftrightarrow\quad
 \mathcal E_{696}\nvdash N(d,d)=d.
\]
Consequently, with the same explicitly retained raw-EPSC calibration,
\[
 \mathrm{ZFC}\vdash
 \bigl(\mathcal E_{692}\nvdash N(d,d)=d
       \ \longleftrightarrow\ \operatorname{Con}(\mathrm{ZFC})\bigr).
\]

\begin{proof}
For the lower direction, take a countable separating
$\mathcal E_{692}$-model and use the same Boolean homomorphism to
two values as in the extraction lemma.  Set
$\nu_0=d$ and $\nu_{i+1}=Q(\nu_i)$.
Equations~\eqref{eq:704-head} and~\eqref{eq:692-tail} recover every
coordinate of a finite stack by induction on $i$.  Represent atomic
membership by $A(\nu_i,\nu_j)$, NAND by $P$, and ordinary existential
quantification by $Q(P(a,a))$.  This still represents every finite
equality-free pure-membership formula.  A selector is an argument of
the atomic constructor; it need not have any particular class truth
value when interpreted as a class in its own right.
Construct the class code before its descriptor, with the same unused
coordinate 2.  The closure equation is unchanged.  The earlier
Boolean extraction therefore gives a countable raw-EPSC structure.

For the upper direction, use precisely the previous free carrier,
class decoding, weighted raw-value recursion, and final Boolean
transport, except for this definition:
\[
 \operatorname{idx}(Q(i))=\operatorname{idx}(i)+1,\qquad
 \operatorname{idx}(i)=0
       \quad\text{if $i$ has no outer $Q$ constructor}.
\]
This is structural recursion on finite trees.  Define $V(s,i)$
to return stream coordinate $\operatorname{idx}(i)$, as before.
The two new native stack equations hold for all $s,x,i$, including
noncanonical indices and tails.  The atomic class $A(i,j)$ reads
the newly defined indices; the class $Q(a)$ retains exactly its
$\exists\neg$ interpretation.

Descriptor descent is unchanged: arbitrarily deep selectors choose
only an input coordinate, a bound coordinate, a proper-subtree
parameter of the descriptor, or the default $d$.  The definition of
$\operatorname{idx}$ inspects syntax and never queries a raw value.
Thus the weight inequalities~\eqref{eq:704-weight} still make the
raw assignments to $Q(b),P(b,z),A(b,z)$ well founded.  The
$W(x,d)$ member-choice slice remains free because class $d$ is false.
The previous proof then verifies every remaining equation and the
separation of the endpoints.

We have obtained a countable raw-EPSC structure from every separating
$\mathcal E_{692}$-model and a separating $\mathcal E_{692}$-model
from every countable raw-EPSC structure.  The earlier argument gives
the same two implications for $\mathcal E_{696}$.  Apply ground-term
equational completeness in both directions to obtain the stated
relative theorem.  No replacement of a native equality by equality
of raw projections occurs in this argument.
\end{proof}

\paragraph{Checking the apparent collision at $Q(d)$.}
In the constructed upper model, $Q(d)$ serves three compatible roles.
As a selector it denotes index 1.  As a class it is true, since its
body $d$ is false and the domain is nonempty.  As a raw object,
$\mu(Q(d))$ is the chosen hull for the constantly false descriptor
relation $R_d$, so hull support makes that raw object empty.
These are different projections of one finite tree.  None of the
equations identifies its class truth value, its selector index, and
its raw-membership value.  In particular, class truth does not
require the tree $Q(d)$ itself to be the Boolean top $d$; the
predicate value $p_\varepsilon(s,Q(d))$ is the top.  This distinction
also explains why no sharing of $V$ with the independent witness
operation $W$ is being asserted.

\paragraph{Complete charge and transport format.}
The eight equation-node counts, totaling 156, are
\[
 14,6,6,9,12,12,17,80.
\]
The codec, local-variable
capacity, opcode width, and endpoints are unchanged:
\[
 \underbrace{8+3\cdot9+8}_{43\text{ header bits}}
 +\underbrace{8+1+4\cdot156}_{633\text{ equation bits}}
 +\underbrace{4\cdot4}_{16\text{ endpoint bits}}
 =\boxed{692\text{ logical bits}}.
\]
The bit-character SHA-256, excluding the final newline, is
\begin{center}\scriptsize\ttfamily
477e5b5de0876ec2d784ca958540117bb2aef9121eaee670925ab2ee8e1455e7.
\end{center}
\texttt{E692.bits} stores exactly these 692 logical bit characters.
\texttt{E692.bin} packs them most-significant bit first into 87 bytes
and appends four zero padding bits.  The supplied
\texttt{E692.bitlength} records 692.  A transport adapter must crop
to that length before calling the strict packet decoder; the padding
is not part of the equational packet.  No new opcode or custom
compression convention is used to obtain the four-bit reduction.

The builder reconstructs the incoming E713 packet, emits and
roundtrips all three new packets, and rejects every strict prefix
and two malformed controls for each.  It also checks the Boolean
identities on all 256 predicate assignments and the new native stack
equations in 14,400 finite-tree regressions.  Those regressions detect
transcription errors in the index convention; the proof above covers
all finite trees and all countable raw-EPSC structures.  The exact
692-bit consistency problem remains to be compiled into a fully
charged binary blank-tape machine.
